\section{Proposed Methodology}
\label{sec:method}

\subsection{Problem formulation and MTN backbone}
Given an observed pedestrian bounding-box trajectory $\mathbf{X}_{obs}\in\mathbb{R}^{B\times15\times4}$, the objective is to predict future bounding boxes $\widehat{\mathbf{X}}_{fut}\in\mathbb{R}^{B\times45\times4}$. The MTN backbone embeds the observed trajectory and adds positional encoding, producing a sequence of width $32$, i.e., $[B,15,32]$. Three MTN encoder layers generate trajectory memory $\mathbf{M}\in\mathbb{R}^{B\times15\times32}$. The original flow-related decoder paths are retained; the decoder output is mapped to four bounding-box coordinates and combined with the final observed-box residual. The dataset protocol and baseline implementation materials follow the official IDD-PeD release \cite{iddped2025,iddped_repo}.

\subsection{Native IDD-PeD context branch}
\label{sec:context}
The first contribution is a context branch built only from six verified native IDD-PeD sources: \emph{intersection\_type}, \emph{signalized\_type}, \emph{road\_type}, \emph{location\_type}, \emph{motion\_direction}, and observed \emph{occlusion}. Day/night and illumination are not included because they were not verified as inputs of the implemented final context branch.

Each source is represented by a 16-dimensional embedding. The five static labels use independent embedding lookups; the observed occlusion history is embedded and mean-pooled over observed steps. The six vectors are concatenated to a 96-dimensional context feature, then projected to the MTN width of 32. The original context-admission gate receives the context representation together with the final observed bounding box, and yields a gated context token $\mathbf{c}\in\mathbb{R}^{B\times1\times32}$. The token is appended \emph{after} the trajectory encoder, yielding augmented decoder memory $[B,16,32]$. This preserves the baseline trajectory memory while allowing the decoder's existing source-attention path to condition on scene context.

\subsection{C4: cross-layer encoder self-attention weight sharing}
\label{sec:c4}
The second contribution is selective cross-layer sharing in the trajectory encoder. C4 aliases only the Q/K/V/O projections of \texttt{EncoderLayer.self\_attn} across encoder layers 0--2. The encoder's \texttt{src\_attn}, \texttt{cross\_attn}, feed-forward networks, normalisation layers, decoder modules, context branch, and output head remain independent. Therefore, C4 is not full layer sharing.

This choice reduces the number of unique parameters while preserving the three sequential encoder-layer executions. It must consequently be interpreted as a parameter-footprint technique, not as a FLOP-reduction claim. The measured CPU latency is reported separately in Section~\ref{sec:results}.

\begin{figure*}[t]
\centering
\resizebox{\textwidth}{!}{% Colourful paper-style overview. Text is intentionally short; detail belongs in the caption.
\begin{tikzpicture}[font=\small,
 t/.style={draw=blue,fill=blue!18,rounded corners=3pt,minimum height=12mm,text width=27mm,align=center},
 c/.style={draw=teal,fill=teal!17,rounded corners=3pt,minimum height=12mm,text width=27mm,align=center},
 g/.style={draw=orange,fill=orange!20,rounded corners=3pt,minimum height=12mm,text width=27mm,align=center},
 o/.style={draw=purple,fill=purple!14,rounded corners=3pt,minimum height=12mm,text width=27mm,align=center},
 a/.style={-Latex,very thick,draw=blue},ca/.style={-Latex,very thick,draw=teal},da/.style={-Latex,dashed,very thick,draw=teal}]
\node[t] (x) {Observed boxes\\$[B,15,4]$};
\node[t,right=10mm of x] (e) {Trajectory\\embedding\\$[B,15,32]$};
\node[t,right=10mm of e] (enc) {MTN Encoder\\$\times 3$};
\node[t,right=10mm of enc] (m) {Trajectory memory\\$[B,15,32]$};
\node[o,right=12mm of m] (aug) {Append context\\token\\$[B,16,32]$};
\node[t,right=10mm of aug] (dec) {MTN Decoder\\$\times3$};
\node[t,right=10mm of dec] (y) {Future boxes\\$[B,45,4]$};
\draw[a] (x)--(e);\draw[a](e)--(enc);\draw[a](enc)--(m);\draw[a](m)--(aug);\draw[a](aug)--(dec);\draw[a](dec)--(y);
\node[c,below=17mm of enc] (share) {C4 shared\\Q / K / V / O};
\draw[da](share)--(enc);
\node[c,below=17mm of x] (scene) {Six IDD-PeD\\context sources};
\node[c,right=10mm of scene] (ctx) {Embeddings +\\projection\\$[B,32]$};
\node[g,right=10mm of ctx] (gate) {Admission\\gate};
\node[o,right=10mm of gate] (token) {Context token\\$[B,1,32]$};
\draw[ca](scene)--(ctx);\draw[ca](ctx)--(gate);\draw[ca](gate)--(token);\draw[ca](token.east)-| (aug.south);
\node[font=\scriptsize,text=blue,below=7mm of ctx] {\textcolor{blue}{Solid: data flow} \qquad \textcolor{teal}{Dashed: shared weights only}};
\end{tikzpicture}
}
\caption{Colour-coded overview of the proposed C4 Context-Aware MTN. Blue: trajectory pathway; teal: native context pathway; orange: admission decision; purple: memory augmentation. Dashed teal denotes shared weights only.}
\label{fig:c4-overview}
\end{figure*}

\begin{figure*}[t]
\centering
\resizebox{\textwidth}{!}{% Context branch as a compact, colourful conceptual mechanism.
\begin{tikzpicture}[font=\small,
 s/.style={draw=teal!70!black,fill=teal!15,rounded corners=2pt,minimum height=9mm,text width=25mm,align=center},
 p/.style={draw=blue!70!black,fill=blue!15,rounded corners=3pt,minimum height=15mm,text width=37mm,align=center},
 g/.style={draw=orange!80!black,fill=orange!20,rounded corners=3pt,minimum height=15mm,text width=40mm,align=center},
 tok/.style={draw=purple!80!black,fill=purple!15,rounded corners=3pt,minimum height=15mm,text width=34mm,align=center},
 a/.style={-Latex,very thick,draw=teal}]
\node[s] (a1) {intersection};\node[s,below=5mm of a1](a2){signalized};\node[s,below=5mm of a2](a3){road};
\node[s,right=6mm of a1](a4){location};\node[s,below=5mm of a4](a5){motion direction};\node[s,below=5mm of a5](a6){occlusion history};
\node[draw=teal,rounded corners=3pt,inner sep=5mm,fit=(a1)(a6),label={[font=\bfseries,text=teal]above:Native context}] (src) {};
\node[p,right=18mm of src] (emb) {Six 16-D\\embeddings};
\node[p,right=14mm of emb] (cat) {Concatenate\\$[B,96]$};
\node[p,right=14mm of cat] (proj) {Context\\projection\\$[B,32]$};
\node[g,right=14mm of proj] (gate) {Gate input\\context + final\\observed box\\$36\rightarrow16\rightarrow2$};
\node[tok,right=14mm of gate] (tok) {Gated context\\token\\$[B,1,32]$};
\draw[a](src)--(emb);\draw[a](emb)--(cat);\draw[a](cat)--(proj);\draw[a](proj)--(gate);\draw[a](gate)--(tok);
\node[font=\small,text=orange!80!black,below=17mm of gate] {The gate decides whether the context token is admitted to decoder memory.};
\end{tikzpicture}
}
\caption{Context-token generation and admission gate. Each native context source is embedded before projection to the MTN width; the gate receives global context and the final observed box.}
\label{fig:c4-context-gate}
\end{figure*}

\begin{figure*}[t]
\centering
\resizebox{\textwidth}{!}{% C4 mechanism: concise and visually close to a method-paper subfigure.
\begin{tikzpicture}[font=\small,
 attn/.style={draw=blue!80!black,fill=blue!15,rounded corners=3pt,minimum height=13mm,text width=39mm,align=center},
 other/.style={draw=gray!80,fill=gray!12,rounded corners=3pt,minimum height=13mm,text width=39mm,align=center},
 bank/.style={draw=teal!80!black,fill=teal!20,rounded corners=3pt,minimum height=15mm,text width=58mm,align=center},
 a/.style={-Latex,very thick,draw=blue},d/.style={-Latex,dashed,very thick,draw=teal}]
\node[attn] (s0) {Layer 0\\self-attention};\node[other,right=8mm of s0](r0){Layer 0\\other modules};
\node[attn,below=16mm of s0](s1){Layer 1\\self-attention};\node[other,right=8mm of s1](r1){Layer 1\\other modules};
\node[attn,below=16mm of s1](s2){Layer 2\\self-attention};\node[other,right=8mm of s2](r2){Layer 2\\other modules};
\draw[a](s0)--(r0);\draw[a](s1)--(r1);\draw[a](s2)--(r2);
\node[bank,right=28mm of r1] (bank) {\textbf{One shared attention bank}\\Q, K, V and O projections};
\draw[d](bank.west)-|(s0.east);\draw[d](bank.west)--(s1.east);\draw[d](bank.west)-|(s2.east);
\node[draw=orange!80!black,fill=orange!14,rounded corners=2pt,align=center,text width=80mm,below=17mm of r1] {\textbf{Not shared:} \texttt{src\_attn}, \texttt{cross\_attn}, FFNs, normalisation, context branch, decoder and output head.};
\end{tikzpicture}
}
\caption{C4 selective cross-layer sharing. Only encoder self-attention Q/K/V/O projections are tied across depth; remaining modules are independent.}
\label{fig:c4-sharing}
\end{figure*}

\begin{figure*}[t]
\centering
\resizebox{\textwidth}{!}{% Token append / decoder detail with a simple token visual.
\begin{tikzpicture}[font=\small,
 mem/.style={draw=blue!80!black,fill=blue!15,rounded corners=3pt,minimum height=17mm,text width=42mm,align=center},
 tok/.style={draw=teal!80!black,fill=teal!20,rounded corners=3pt,minimum height=17mm,text width=38mm,align=center},
 aug/.style={draw=purple!80!black,fill=purple!15,rounded corners=3pt,minimum height=17mm,text width=45mm,align=center},
 outputbox/.style={draw=orange!80!black,fill=orange!20,rounded corners=3pt,minimum height=17mm,text width=42mm,align=center},
 a/.style={-Latex,very thick,draw=blue},ca/.style={-Latex,very thick,draw=teal}]
\node[mem](memory){15 trajectory tokens\\$[B,15,32]$};
\node[tok,below=21mm of memory](token){1 gated context token\\$[B,1,32]$};
\node[aug,right=28mm of memory,yshift=-11mm](append){Append\\16 decoder-memory tokens\\$[B,16,32]$};
\node[mem,right=25mm of append](decoder){MTN Decoder $\times3$\\flow paths unchanged};
\node[outputbox,right=25mm of decoder](output){Output head + residual\\future boxes\\$[B,45,4]$};
\draw[a](memory)--(append);\draw[ca](token.east)-|(append.south);\draw[a](append)--(decoder);\draw[a](decoder)--(output);
\end{tikzpicture}
}
\caption{Decoder-memory augmentation. The gated context becomes the sixteenth decoder-memory token without replacing the original 15 trajectory tokens.}
\label{fig:c4-memory}
\end{figure*}

\subsection{Novel contributions}
\begin{enumerate}[leftmargin=*,label=\textbf{C\arabic*.}]
\item \textbf{Context-token augmentation of MTN.} The method introduces six verified IDD-PeD context sources as a gated 32-dimensional decoder-memory token without replacing the trajectory pathway.
\item \textbf{Selective, technically scoped encoder sharing.} C4 ties only encoder self-attention Q/K/V/O across depth. This avoids incorrectly claiming that all encoder/decoder components are shared or that FLOPs decrease.
\item \textbf{Matched evidence-driven architecture selection.} B, C4, F, and F+C4 were evaluated under a common small CPU protocol, enabling a factual comparison of accuracy, variation across seeds, parameter count, and latency.
\end{enumerate}

\subsection{Separate Dynamic Multimodal Fusion candidate}
Dynamic Multimodal Fusion (F) was implemented and evaluated separately. It uses six 32-dimensional context tokens for cross-attention and a separate global context vector. The trajectory memory acts as query, while the six context tokens are keys and values. Identity, additive, and cross-attention candidates are computed, then a $64\rightarrow32\rightarrow3$ gate selects one operation per sample using hard Gumbel-Softmax in training and argmax/one-hot selection in evaluation. All three candidates are currently computed before selection; therefore F has no branch-skipping or resource-aware computational-saving claim. F is not included in the C4 proposed architecture in Figure~\ref{fig:c4-overview}.
